\chapter{Conclusion}
\label{ch:conclusion}

\section{Summary of Work}
\label{sec:summary}

FoodLink set out to shorten the time between surplus food being posted and a
trustworthy organisation accepting it, and to see the food through delivery
with a record of every step. By mid-term, the team has built and verified the
core system, meeting Objectives 1--5 and 7. The system comprises a FastAPI
backend (27 routes, six tables) and a React frontend with four role portals.
Donations follow a nine-state lifecycle with conditional writes and an
append-only event ledger. An explainable weighted matcher ranks organisations,
using requirements only to break ties. Role, ownership, lifecycle and
verification controls are enforced, with disclosure tailored to each reader.
Standing requirements, ledger-derived metrics and administrator tools
complete the feature set, and 495 backend and 156 frontend tests pass in CI.
Objective~6 exists as instrumentation awaiting real use, and deployment and a
pilot (Objective~8) have not started. The limitations in
Section~\ref{sec:limitations} remain open, notably the missing portal controls
and the organisation map-pin gap.

\section{Key Findings}
\label{sec:findings}

\textbf{Finding 1:} \emph{a transparent heuristic is sufficient and
verifiable.} Rankings can be recomputed by hand, and tying explanations to the
scoring code exposed a real defect: two criteria that were effectively one.
\textbf{Finding 2:} \emph{most of the difficulty lies in access and integrity,
not features.} The most serious defects were unverified accounts reading donors'
addresses, verification that survived a change of identity, and double
acceptance, all invisible in single-user testing. \textbf{Finding 3:}
\emph{race conditions are real even on SQLite.} Moving each precondition into
the write itself fixed them testably. \textbf{Finding 4:} \emph{privacy needs
reader-specific answers.} Distances and scores can reveal the locations they
come from, so what a record says must depend on who reads it.

\section{Future Work and Recommendations}
\label{sec:future}

None of the following is implemented. Recommended order:
\begin{enumerate}
  \item close the remaining audit items: an exception handler, profile bounds,
    a guarded expiry sweep, the deep-link fix, removal of the mobile residue,
    and frontend tests;
  \item complete the portal with a map-pin field for organisations and
    controls for cancel, release and account management;
  \item deploy with a secret key, a scheduled sweep and, if needed, PostgreSQL;
    then run a campus pilot to measure time-to-claim, rescue rate and handover
    time;
  \item add notifications, token revocation, security headers and an
    administrator audit log;
  \item add a map-based pin picker, object storage for photos, pagination and
    database-side aggregation;
  \item extend matching with structured food categories and road distance.
    Weight tuning should wait for outcome data, and any learned ranker would sit
    behind \code{score_pair}, must remain explainable, and must show measured
    improvement;
  \item prepare legally: a privacy notice and consent, data-access and deletion
    requests, retention rules, and terms of use.
\end{enumerate}

\section{Final Remarks}
\label{sec:remarks}

FoodLink has moved from a proposal to a working, tested system whose rules can
be read, whose rankings can be checked by hand, and whose history cannot be
rewritten by its users. The most valuable lessons came from treating the system
adversarially. What remains before the final evaluation is to put it in front of
real users and let the event ledger measure whether surplus food reaches people
before it expires.
